\section{Finite quotients and approximation}\label{sec:foundations}

The first objective is to isolate what local approximation does not already
prove.  Every noncentral normal subgroup has finite index, but its closure
at the anisotropic finite places need not a priori recover the subgroup.
We express this possible discrepancy as a finite quotient.  We then
develop two tools for eliminating it: an approximation procedure that
retains the abstract rational coset, and a Hasse principle with weak
approximation for a particular norm torus.

\subsection{Conventions and established inputs}

Throughout, \(k\) is a number field, \(V_k\) is its set of places, and
\(G\) is an absolutely almost simple simply connected \(k\)-group.
Write
\[
 A=\{v\in V_k:\ v\text{ is nonarchimedean and }
                    \rank_{k_v}G=0\},\qquad
 G_A=\prod_{v\in A}G(k_v).
\]
The diagonal homomorphism to this product is denoted by \(\delta\).
Thus \(G_A=H_A\) and \(\delta=\delta_A\) in the notation of the introduction.
An empty product is the trivial group.  The set \(A\) is finite: outside
finitely many places the group has a reductive model and is unramified
and quasi-split, hence has positive relative rank.
At a nonarchimedean place an absolutely simple anisotropic group has
type \(A\).

For a finite set of places \(B\), a subscript \(B\) on a rational element
denotes its tuple of local images.  Products over finitely many places
have their ordinary product topology.  In an adelic product we use the
restricted product with respect to fixed integral compact open subgroups
outside a finite set.  Weak approximation means density at finitely many
completions; strong approximation means density in the specified
restricted adelic product.

For an integer \(e\geq 1\), the notation
\[
                 G(k)^e=\langle g^e:g\in G(k)\rangle
\]
always denotes the abstract subgroup generated by the powers.  No closure
is included in this definition.  Degrees of central simple algebras are
reduced degrees; for example, \(\SL_1(M_r(\Delta))\) has type
\(A_{r\deg\Delta-1}\).  Restrictions of scalars will be treated over
their defining extension fields when applying an induction or an
approximation theorem.

We use the strong approximation theorem for simply connected absolutely
almost simple groups: if \(G(k_w)\) is noncompact, then \(G(k)\) is dense
in the adelic product away from \(w\)
\cite[Theorem 7.12]{PRbook}.  We also use the lattice theorem
for semisimple \(S\)-arithmetic groups
\cite[Theorem~5.7(1)]{PRbook}, the normal subgroup theorem for
irreducible \(S\)-arithmetic lattices of total rank at least two, and the
local Kneser--Tits and simplicity theorems.  In the last assertion, for
a nonarchimedean \(v\) with positive rank, \(G(k_v)\) is perfect and
\(G(k_v)/Z(G(k_v))\) is abstractly simple.
These results are used in their number-field forms
\cite[Chapters VII and IX]{PRbook}\cite[Chapter VIII, Theorem A]{Margulis}.

We refer to the assertion of Theorem~\ref{thm:main} for a particular
group over a particular number field as \emph{(MP)}.  We require the
following established cases.

\begin{theorem}[Known cases of (MP)]\label{thm:known-mp}
The assertion (MP) holds except possibly for anisotropic groups of
outer type \(A\), trialitarian type \(D_4\), and type \(E_6\).
In particular it holds for every inner form of \(\SL_n\), over every
number field.
\end{theorem}

The formulation, including the results for multiplicative groups of
division algebras that enter the inner type \(A\) case, is recalled in
\cite[Sections 3.4--3.5]{PRsurvey}. For the isotropic completion see
\cite[Theorem~8.1]{Gille2009}; for inner type \(A\) see
\cite{RapinchukPotapchik1996,Segev1999,SegevSeitz2002}, with the subsequent
strengthening in \cite{RapinchukSegevSeitz2002}.
Thus, when we address one of the remaining types, we may assume that
the group is anisotropic over its ground field.

We will also use the following form of the Hasse principle.  A torsor
under a \(k\)-group \(H\) means a principal homogeneous space for \(H\)
over \(k\).

\begin{theorem}[Simply connected Hasse principle]
\label{thm:foundations-hasse}
Let \(H\) be a simply connected semisimple group over a number field
\(k\).  Then \(H^1(k_v,H)=1\) at every nonarchimedean place.  A torsor
\(X\) under \(H\) that has a point over every completion of \(k\) has
a \(k\)-point.  Moreover, \(X(k)\) is dense in
\(\prod_{v\in B}X(k_v)\) for every finite set \(B\) of places.
\end{theorem}

The local and global vanishing assertions are the Kneser--Harder--Chernousov
Hasse principle; we use the standard number-field statements in
\cite[Theorems~6.4 and~6.6]{PRbook}.  For the last assertion, a rational point identifies
\(X\) with \(H\).  Decompose \(H\) as a product of restrictions of scalars
of simply connected absolutely almost simple groups.  For each factor,
choose an auxiliary split finite place outside the given approximation
set and apply strong approximation away from it.  This gives the stated
weak approximation, also for products and restrictions of scalars.

We use the usual local and global theory of central simple algebras
and involutions, including local Brauer invariants, the embedding
criterion for fields, the Hasse--Schilling norm theorem, and the
classification of hermitian forms
\cite{PRbook,Reiner,KMRT}.  When an etale algebra is a product of fields,
norm and local arguments are understood factor by factor.  All fields
have characteristic zero unless they are explicitly residue fields.
In root computations the roots of a simply laced system have squared
length two.  A superscript \(1\) on a multiplicative group denotes
the kernel of the indicated norm.

\subsection{The finite quotient not detected by closure}

\begin{proposition}[Finite-defect reduction]\label{prop:finite-defect}
Every abstract noncentral normal subgroup of \(G(k)\) has finite index.
Fix \(e\geq1\), and put
\[
 R=G(k)^e,\qquad
 P_0=\overline{\delta(R)}\subset G_A,\qquad
 V_0=\delta^{-1}(P_0),\qquad V=V_0/R.
\]
Then \(R\) has finite index, \(P_0\) is open and normal in \(G_A\),
and \(V\) is finite.  For every positive-rank nonarchimedean place
\(w\), the closure of \(R\) in the adelic product away from \(w\)
is exactly the inverse image of \(P_0\).

To prove (MP) for \(G\), it suffices to show \(V=1\) for every \(e\).
If \(\gamma\in V_0\), representatives of \(\gamma R\) can be chosen
arbitrarily close to \(1\) at any prescribed finite set of places.
Away from \(A\), these near-identity conditions may be replaced by
arbitrary nonempty local open conditions.  Finitely many nonempty
Zariski open conditions may be imposed simultaneously.
\end{proposition}

\begin{proof}
Let \(N\lhd G(k)\) be noncentral, and choose a noncentral \(n\in N\).
Take a finite set \(T\) containing the archimedean places, \(A\),
the denominator places of \(n\), and auxiliary split finite places
whose total local rank is at least two.  For fixed integral compact
opens outside \(T\), set
\[
 K^T=\prod_{v\notin T}G(\mathcal O_v),\qquad
 \Gamma=G(k)\cap K^T,
\]
where changing finitely many integral models changes neither argument.
The group \(\Gamma\) is the corresponding \(T\)-arithmetic group and
contains \(n\).

The lattice theorem \cite[Theorem~5.7(1)]{PRbook} makes \(\Gamma\) a lattice in the product of the
local groups at \(T\).  Remove its compact factors.  This preserves
discreteness: the inverse image of a compact subset of the remaining
product is compact after the compact factors are restored, and hence
meets \(\Gamma\) in finitely many points.  It also preserves finite
covolume.  The rational embedding remains faithful.  Strong
approximation away from each noncompact factor supplies the projection
density required for irreducibility.  The \(S\)-arithmetic normal
subgroup theorem now applies to \(\Gamma\cap N\), which is normal and
noncentral, and yields
\[
                         [\Gamma:\Gamma\cap N]<\infty .
\]
If one uses the formulation with a finite-normal-subgroup alternative,
that alternative is central: the centralizer of a finite normal
subgroup contains a finite-index Zariski-dense subgroup of \(\Gamma\),
and hence is all of the connected group \(G\).

Strong approximation, omitting a positive-rank place in \(T\), shows
that \(\Gamma\) is dense in \(K^T\).  Thus the closure of
\(\Gamma\cap N\) has finite index in \(K^T\).  It is closed, so it
is open.  It follows that the closure \(C\) of \(N\) in the restricted
product away from \(T\) is open.  It is normal because \(N\) is
normal and \(G(k)\) is dense there.

Every place outside \(T\) has positive rank.  For such a place \(v\),
the intersection \(C\cap G(k_v)\), with the local group supported at
that single place, is an open normal subgroup.  It cannot be central,
since the center is finite.  Local simplicity modulo center and
perfectness give \(C\cap G(k_v)=G(k_v)\).
Finite-support products are dense in the restricted product, so
\(C\) is the whole restricted product.  Given \(g\in G(k)\), we can
therefore choose \(a\in N\) with \(ag\in K^T\), and hence \(ag\in\Gamma\).
The map from \(\Gamma/(\Gamma\cap N)\) to \(G(k)/N\) is surjective.
This proves the finite-index assertion.

The power map has differential \(e\) times the identity at \(1\);
it is therefore dominant in characteristic zero.  Rational points
of a connected linear algebraic group over \(k\) are Zariski dense.
Consequently \(R\) is Zariski dense and noncentral, so the preceding
assertion applies to \(R\).
Density of \(G(k)\) in \(G_A\) follows from strong approximation away
from an auxiliary positive-rank finite place.  Since \(R\) has finite
index, its closure \(P_0\) has finite index in \(G_A\), is normal, and
is open.

Now omit any positive-rank finite place \(w\).  The closure \(C_w\)
of \(R\) in the remaining adelic product is again closed, normal,
and of finite index: finitely many translates of it already contain
the dense subgroup \(G(k)\).  Thus it is open.  The local argument
just given shows that \(C_w\) contains every positive-rank
nonarchimedean factor.  It contains every archimedean factor as well:
the group of real points of a simply connected semisimple algebraic
group is connected, and so is the group of complex points.  Their
maps into the finite discrete quotient by \(C_w\) are therefore
trivial.  Taking closures of finite-support products proves that
\(C_w\) contains the kernel of the projection to \(G_A\).
Its image there is closed, contains \(\delta R\), and is contained
in \(P_0\); hence it is exactly \(P_0\).  This proves the closure
description.

For an original noncentral normal subgroup \(N\), the finite quotient
\(G(k)/N\) has some exponent \(e\), so \(R\subset N\).  If \(V=1\),
then \(R=\delta^{-1}(P_0)\), and density gives an isomorphism
\[
                         G(k)/R \simeq G_A/P_0 .
\]
The normal subgroup \(N/R\) corresponds to a normal subgroup of this
finite quotient.  Its inverse image in \(G_A\) is an open normal
subgroup \(W\), and \(N=\delta^{-1}(W)\).

Finally, for \(\gamma\in V_0\), the closure of \(\gamma R\) away
from \(w\) is the same inverse image of \(P_0\), since
\(\gamma_A P_0=P_0\).  It contains the identity at the places in \(A\)
and all tuples at other places.  Choose \(w\) outside any prescribed
finite approximation set to obtain the asserted local choices.
Nonempty local open subsets are Zariski dense in the smooth connected
group; intersecting one such opening with any additional nonempty
Zariski open subset preserves nonemptiness.  This supplies the last
assertion.
\end{proof}

The following consequence records exactly how an already established
case of (MP) will be used for subgroups.  In particular it does not
require a finite homomorphism to be continuous in advance.

\begin{lemma}[Restriction to a group satisfying (MP)]
\label{lem:restriction-mp}
Let \(\psi:G(k)\to F_0\) be a homomorphism to a finite group of
exponent dividing \(e\).  Let \(f:H\to G\) be a \(k\)-homomorphism,
where \(H\) is a product of restrictions of scalars of simply connected
absolutely almost simple groups for which (MP) is known over their
respective number fields.  The homomorphism \(\psi\circ f\) on \(H(k)\)
extends continuously to the product of the anisotropic finite local
groups of those factors.  Its image still has exponent dividing \(e\).
Consequently \(\psi(f(h))=1\) if \(h\) is an \(e\)-th power at every
one of these local factors.
\end{lemma}

\begin{proof}
For one absolutely almost simple factor, the kernel of the restricted
finite homomorphism has finite index and is Zariski dense, so it is
not central.  By (MP) it is the inverse image of an open normal
subgroup \(W_H\) of the anisotropic local product.  Density identifies
the rational quotient with the finite quotient by \(W_H\).
The homomorphism therefore extends through that quotient.
For a product, extend the maps factor by factor.  Their images
commute, as they do on rational points, so the extensions combine.
Apply this argument over the extension field for each restriction
of scalars.  The extended image is unchanged, and every local
\(e\)-th power maps to \(1\).
\end{proof}

\subsection{An affine avoidance lemma}

The closure calculation permits finitely many local conditions.
The next argument supplies an additional condition at almost all
finite places.  It will be applied to an affine space of parameters,
not to integral points of a torus.
Codimension-two avoidance in affine space is an instance of arithmetic
purity of strong approximation; see \cite[Proposition~3.6]{CaoXu2018}
and \cite[arXiv version, Lemma~1.1]{Wei2019}. We give the fixed-direction version with
an unbounded parameter at the omitted place needed below.

\begin{lemma}[Avoidance along a line]\label{lem:foundations-affine-sieve}
Let \(E\) be an \(N\)-dimensional \(k\)-vector space, and let
\(Z\subset E\) be closed of geometric codimension at least two.
Choose \(d\in E(k)\setminus\{0\}\) whose point at infinity is outside
the projective closure of \(Z\), and let
\(\pi:E\to E/kd\) be the quotient map.  Fix a nonarchimedean place
\(v_0\), linear coordinates, and a rational linear section of \(\pi\).
Outside a finite set of model exceptions, require that these data
give an integral direct sum \(E=kd\oplus E'\), with \(d\) an integral
basis for the first summand, that \(\pi|_Z\) remains finite on
reduction, and that every geometric fiber has cardinality at most
\(D\).

Let \(\Sigma\) be a finite set containing \(v_0\), the archimedean
places, those model exceptions, and all places with residue
cardinality at most \(D\).  For every
\(v\in\Sigma\setminus\{v_0\}\), prescribe a nonempty open subset
\(\mathcal O_v\subset E(k_v)\).
There are a vector \(y\in E(k)\) and elements \(l\in k\) with
\(|l|_{v_0}\) arbitrarily large such that \(a=y+ld\) belongs to
\(\mathcal O_v\) at these prescribed places, is integral outside
\(\Sigma\), and its reduction avoids \(Z\) at every finite place
outside \(\Sigma\).
\end{lemma}

\begin{proof}
The condition on the point at infinity makes \(\pi|_Z\) finite.
For example, extend the projection to projective space: the only
indeterminacy is the chosen point at infinity, which is disjoint
from the projective closure of \(Z\).  Over an affine point of the
target, the only possible point at infinity on its fiber line is
that same point.  Thus the restricted morphism is proper and has
finite fibers, and is finite.  Its image \(Y\subset E/kd\) is closed
and proper, since
\(\dim Z\leq N-2<\dim(E/kd)\).
Use the fixed linear coordinates and integral models; their
exceptional places already belong to \(\Sigma\).

Linear projection is open at each completion.  Additive strong
approximation away from \(v_0\) gives
\[
 q\in (E/kd)(k)\setminus Y(k)
\]
that is integral outside \(\Sigma\) and belongs to
\(\pi(\mathcal O_v)\) at every prescribed place.
To ensure \(q\notin Y\), first shrink one of the prescribed
nonempty openings in the quotient to miss \(Y\); a proper algebraic
subset has empty interior.  Choose a rational linear lift \(y\) of
\(q\), integral outside \(\Sigma\).

Because \(q\notin Y\), its reduction belongs to the reduction of
\(Y\) at only finitely many finite places outside \(\Sigma\).
Indeed a polynomial vanishing on \(Y\), but not at \(q\), is a
nonzero integral value away from finitely many places.
At each of the remaining exceptional places, the forbidden points
in the fiber line number at most \(D\).  The residue field has
more than \(D\) elements, so an integral value of \(l\) avoiding
them exists.  This is a nonempty local congruence condition on \(l\).
At a prescribed place \(v\), the condition \(q\in\pi(\mathcal O_v)\)
likewise gives a nonempty local open set of \(l\) with
\(y+ld\in\mathcal O_v\).

Apply additive strong approximation once more, now to \(l\), imposing
these finitely many conditions and integrality at all other finite
places outside \(\Sigma\).  The local openings may be chosen
relatively compact away from \(v_0\).  There are infinitely many
such rational \(l\), by density in the adeles away from \(v_0\).
They cannot all remain bounded at \(v_0\), since the diagonal
copy of \(k\) is discrete in its full adele ring.  We may therefore
choose \(|l|_{v_0}\) arbitrarily large.  This gives all the assertions.
\end{proof}

\subsection{Approximation by abstract power factors}

We now use anisotropy to make the unrestricted direction in
Lemma~\ref{lem:foundations-affine-sieve} harmless at \(v_0\).
Assume \(G\) is \(k\)-anisotropic.  Fix a maximal \(k\)-torus
\(T_0\subset G\), a finite Galois extension \(P/k\) splitting \(T_0\)
and any fixed auxiliary data, and a regular cocharacter
\(\lambda:\bbG_{m,P}\to T_{0,P}\).
The map
\begin{equation}\label{eq:foundations-eta}
 \eta:\Res_{P/k}\bbG_m\longrightarrow T_0,\qquad
 \eta(a)=N_{P/k}\bigl(\lambda(a)\bigr)
\end{equation}
is defined over \(k\).  Its restriction to the diagonal \(\bbG_m\)
has cocharacter \(\sum_{\sigma\in\Gal(P/k)}\lambda^\sigma\).
This is a \(k\)-cocharacter and hence is zero by anisotropy.
In particular
\begin{equation}\label{eq:foundations-scalar-invariance}
                         \eta(la)=\eta(a)\quad(l\in k^\times).
\end{equation}

For \(x\in G(k)\), or for \(x\) in a fixed rational algebraic
coset \(X=x_*G\) inside a larger linear algebraic group, consider
\begin{equation}\label{eq:foundations-power-map}
 \Phi(a_1,\ldots,a_r)
       =x\prod_{j=1}^r k_j\eta(a_j)^e k_j^{-1},
       \qquad k_j\in G(k),\quad a_j\in P^\times .
\end{equation}
Its values lie in the same right coset \(xR\), with
\(R=G(k)^e\).  Its parameter domain is the open torus
\((\Res_{P/k}\bbG_m)^r\) in the affine \(k\)-space
\(E=(\Res_{P/k}\bbA^1)^r\).
Over \(P\), its boundary consists of the hyperplanes
\(a_{j,\sigma}=0\), one for each absolute coordinate.

We call a reduction a \emph{single-pole reduction} if exactly one
of these absolute coordinates vanishes.  Away from the other
hyperplanes, Equation~\ref{eq:foundations-power-map} then takes the
form
\begin{equation}\label{eq:foundations-single-pole}
             \Phi=g_L\,\lambda^\sigma(t)^e\,g_R,
                    \qquad t=a_{j,\sigma},
\end{equation}
where the side factors are regular group- or coset-valued functions.
At a good finite place with this reduction, Frobenius fixes the unique
vanishing coordinate.  Its action on the embeddings of the Galois
extension \(P/k\) is regular.  Thus the place splits completely in
\(P\), \(t\) has positive valuation, and the side factors are integral.
The local analysis of Equation~\ref{eq:foundations-single-pole} may
therefore be carried out in split coordinates.

Here is the precise criterion used in later sections.  Its hypothesis
about the single-pole configurations is substantive and will be
verified in each application.  A consecutive list of factors in
Equation~\ref{eq:foundations-power-map} is called a \emph{basic list}
if its product map is dominant and has surjective differential
somewhere.

\begin{proposition}[Power-factor approximation]
\label{prop:power-approximation}
Let \(G\) be \(k\)-anisotropic, with \(P\), \(\eta\), \(R\), and
\(X=x_*G\) as above.  Fix \(\gamma\in G(k)\), the rational coset
\(\mathcal C=x_*\gamma R\subset X(k)\), a
geometrically codimension-at-least-two closed subset \(D\subset X\),
and a finite set \(B\) containing the archimedean places, \(A\), and
the exceptions for fixed integral models and sufficiently small
residue fields.  Include in \(B\) a positive-rank finite place \(v_0\).
At \(B\) prescribe nonempty local open conditions compatible with
the closure of \(\mathcal C\) in
Proposition~\ref{prop:finite-defect}.

At each finite place outside \(B\), prescribe an allowed set
\(\mathcal U_v\subset X(k_v)\).  Assume:
\begin{enumerate}
 \item every integral point whose reduction avoids \(D\) belongs
       to \(\mathcal U_v\);
 \item \(\mathcal U_v\) contains a nonempty local open set;
 \item single-pole configurations are tested by finitely many
       algebraic nonvanishing conditions over \(P\), independent
       of \(v\).  Their nonemptiness is compatible with the following
       order of choices.  First choose at least two disjoint basic
       lists, with their conjugator tuple in a nonempty Zariski
       open subset of its product of copies of \(G\).  For each
       such tuple there is a nonempty Zariski open subset of \(X\)
       in which to choose \(x\).  After \(x\) is fixed, for every
       finite number of additional factors there is a nonempty
       Zariski open subset of the corresponding product of copies
       of \(G\) in which to choose their conjugators.  For every
       tuple allowed by these successive choices, the tests
       restrict to a nonempty open subset of each absolute
       boundary hyperplane away from the others.  They guarantee
       membership in \(\mathcal U_v\) at a single-pole reduction,
       outside finitely many model exceptions for the fixed data.
\end{enumerate}
In the third condition the tests are imposed with all their Galois
conjugates.  Appending additional factors is required to preserve
the tests on an existing hyperplane when the new parameters are
set to \(1\).  Thus their insertion does not remove the previously
available open subsets.  The additional factors permit the local
moves required by the second condition.

Then there is \(x'\in\mathcal C\) satisfying the prescribed conditions
at \(B\) and belonging to \(\mathcal U_v\) for every finite
\(v\notin B\).  Finitely many nonempty Zariski open conditions on
the initial representative and conjugators can be included in their
choice, provided they are compatible with the third condition.
\end{proposition}

\begin{proof}
We give the construction in an order that keeps the finite model
exceptions separate from the infinitely many avoidance conditions.
When the preliminary choices enlarge \(B\), use at each newly added
place a nonempty opening inside \(\mathcal U_v\), as supplied by the
second hypothesis.  Such places lie outside \(A\), so these added
conditions are compatible with the coset closure.  Meeting them
preserves the conclusion requested for the original \(B\).

\emph{Dominant parameter lists.}
The adjoint translates of the differential of a nonzero cocharacter
span \(\Lie(G)\): their span is a nonzero adjoint-invariant subspace
of the simple Lie algebra.  Since \(G(k)\) is Zariski dense, finitely
many rational conjugators make the product of their \(\eta^e\)-maps
have surjective differential at some point.  Such a product map is
dominant.  Choose at least two disjoint consecutive lists with this
property in the nonempty open set supplied by the third hypothesis.
We may intersect that set with any further prescribed nonempty
Zariski open conditions.  The product of copies of \(G\) is
geometrically irreducible, so such intersections remain nonempty.

On the open parameter torus the total product map is smooth whenever
one of its basic list maps is smooth.  Each list has a proper closed
nonsmooth locus.  Since the two lists have disjoint parameter sets,
their simultaneous nonsmooth locus has codimension at least two.
The inverse image of \(D\) has codimension at least two on the smooth
locus, and any remaining part lies in that same codimension-two
nonsmooth locus.  Therefore its closure in \(E\) has codimension at
least two.  Adding further independent lists does not change this
conclusion.

We also need a residue-field observation about one fixed basic list.
After discarding finitely many model exceptions, its dominance
persists on reduction.  For any integral translating point \(x\),
the inverse image of the dense open \(X\setminus D\) is a nonempty
open in its affine parameter space.  The equations excluded here,
including the norm equations defining the torus boundary, have
bounded degree independently of the translating point.  The
elementary bound for zeros of a nonzero polynomial over a finite
field therefore supplies a rational parameter in this open whenever
the residue field is sufficiently large.  Thus unit parameters in
one basic list, all other parameters set to \(1\), can give a
reduction outside \(D\).  Put these fixed exceptions and small
residue fields into \(B\).

\emph{The initial point and its denominator places.}
Choose \(x\in\mathcal C\) satisfying the conditions at \(B\), using
Proposition~\ref{prop:finite-defect}; to approximate also at \(v_0\),
omit a different positive-rank place.  Any prescribed compatible
Zariski open conditions can be included; in particular, take \(x\)
in the open subset supplied by the third hypothesis after the
basic conjugators have been fixed.  Let \(B_1\) be the finite
set of places outside \(B\) where \(x\) is not integral.
At each such place choose a nonempty open subset of \(\mathcal U_v\).

The local normal subgroup generated by
\(\eta((P\otimes_k k_v)^\times)^e\) is all of \(G(k_v)\).
Indeed this image is Zariski noncentral, as can be seen in the
independent absolute coordinates, and local simplicity modulo center
and perfectness apply.  Every element of that generated subgroup
is a finite product of conjugates of such powers; inverse factors
are again of the same form.  Consequently finitely many extra
lists can move \(x_v\) into the chosen opening at each \(v\in B_1\).
Fix the finite number of factors needed for these moves, padding
the local lists by factors with parameter \(1\) where necessary.
The third hypothesis supplies a nonempty Zariski open subset for
this whole tuple of extra conjugators, with the initial \(x\)
and basic conjugators unchanged.  Each product of the required
local openings is Zariski dense and hence meets this subset.
Choose their conjugators globally by strong approximation away from
\(v_0\), approximating the required local conjugators and remaining
integral outside \(B\cup B_1\).  Additional nonempty Zariski open
conditions are compatible with this choice.
The basic parameters can be taken near \(1\) on \(B\cup B_1\);
the extra parameters can be taken near \(1\) on \(B\).  We now have
nonempty open conditions on the full affine parameter space at
every place of \(B\cup B_1\).

\emph{The closed subset to be avoided.}
Fix the constants just chosen.  In \(E\), let \(Z\) be the union of
the closure of \(\Phi^{-1}(D)\), intersections of distinct boundary
hyperplanes, and the excluded proper closed subsets of each
individual hyperplane from the third hypothesis.
Take all Galois conjugates, so \(Z\) is defined over \(k\).
Every component has codimension at least two: this was proved for
the first part, is immediate for the intersections, and follows
from nonemptiness of the specified open tests for the last part.
For parameters integral at a good place and reducing outside \(Z\),
there are exactly two possibilities.  Either every torus coordinate
is a unit and \(\Phi\) is integral with reduction outside \(D\), or
there is one pole and its prescribed open test holds.  In both cases
\(\Phi\) belongs to \(\mathcal U_v\).

\emph{The direction and the residual exceptions.}
Choose a rational direction \(d=(d_1,\ldots,d_r)\in E(k)\)
whose point at infinity avoids the projective closure of \(Z\).
Choose each \(d_j\) sufficiently close to \(1\) at \(v_0\) so that
\(\Phi(d)\) still satisfies the prescribed condition there.
These requirements are compatible: avoiding a proper algebraic
subset is dense in any local open set.
The projection along \(d\), its linear splitting, the tests defining
\(Z\), and the fiber bound introduce only finitely many further
exceptions.  They are fixed before the quotient point in the sieve
is chosen.

At a new exceptional place outside \(B\cup B_1\), the point \(x\)
and all conjugators are integral.  Use the residue-field observation
for a basic list to choose unit parameters giving reduction outside
\(D\), and take the extra parameters near \(1\).  This provides a
nonempty local opening where \(\Phi\in\mathcal U_v\), even though
the projection or single-pole model might be unusable there.
Any additional small residue fields needed for the fiber bound
are treated in exactly this way; the smaller fields for which the
basic-list observation itself fails already lie in \(B\).
Let \(\Sigma\) contain \(B\cup B_1\) and all these residual
exceptions.

We can now apply Lemma~\ref{lem:foundations-affine-sieve} with the
parameter openings just specified, omitting the condition at
\(v_0\).  It gives parameters \(a=y+ld\), integral outside
\(\Sigma\), reducing outside \(Z\) there, and meeting all the
openings at \(\Sigma\setminus\{v_0\}\).  As
\(|l|_{v_0}\to\infty\), Equation~\ref{eq:foundations-scalar-invariance}
gives
\[
 \eta(a_j)=\eta(d_j+l^{-1}y_j)\longrightarrow\eta(d_j).
\]
For sufficiently large \(|l|_{v_0}\), every \(a_j\) is nonzero and
\(\Phi(a)\) satisfies the required condition at \(v_0\).
It satisfies the conditions at all other exceptional places by
construction, and lies in \(\mathcal U_v\) elsewhere by avoidance
of \(Z\).  Finally \(\Phi(a)\in xR=\mathcal C\), since every factor
in Equation~\ref{eq:foundations-power-map} is an abstract \(e\)-th
power.  This completes the construction.
\end{proof}

The proposition allows poles, but only in a controlled form at places
that split completely in \(P\).  In particular, it makes no assertion
that the torus parameters can be integral units everywhere.
The later applications will specify \(D\), prove the nonemptiness
of every single-pole test, and check its local consequence.

\subsection{Cohomology detected on cyclic subgroups}

Our final preliminary tool solves simultaneous norm equations with
prescribed local approximations.  We first separate the arithmetic
duality input from the finite-group calculation.

For a discrete Galois module \(M\), write
\[
 \Sha^i(k,M)=
 \ker\left(H^i(k,M)\longrightarrow\prod_{v\in V_k}H^i(k_v,M)\right).
\]
The notation \(\Sha_\omega^i(k,M)\) means classes whose localizations
vanish at all but finitely many places.  For a group \(H\), the
degree-one notation has its pointed-set meaning.  Galois cohomology
and its cochains are continuous throughout.
A \emph{lattice} is a finitely generated free abelian group with a
continuous Galois action, which necessarily has finite image.

\begin{lemma}[Cyclic detection]\label{lem:foundations-cyclic-detection}
Let \(M\) be a \(k\)-lattice split by a finite Galois extension
\(P/k\), and put \(\Gamma=\Gal(P/k)\).  Inflation identifies
\(\Sha_\omega^2(k,M)\) with
\[
 \ker\left(
 H^2(\Gamma,M)\longrightarrow
             \prod_{C\subset\Gamma\ {\rm cyclic}}H^2(C,M)
 \right).
\]
If this group is zero and \(T\) is a \(k\)-torus with character
lattice \(M\), every everywhere locally soluble torsor under \(T\)
has a rational point and satisfies weak approximation.
\end{lemma}

\begin{proof}
Over \(P\) the action on \(M\) is trivial.  The exact sequence with
rational coefficients identifies
\[
 H^2(P,M)\simeq H^1(P,M\otimes(\bbQ/\bbZ)).
\]
The right side consists of continuous finite-image characters.
If a class is locally zero almost everywhere, its restriction to
\(P\) is zero by Chebotarev.  Moreover
\(H^1(P,M)=0\), since a continuous homomorphism from a profinite
group to a torsion-free discrete abelian group has finite, hence
trivial, image.  Inflation--restriction therefore identifies such
a class with a unique class in \(H^2(\Gamma,M)\).

At a place unramified in \(P\), the decomposition group is cyclic.
Local inflation is injective by the same vanishing of degree-one
cohomology over the splitting field.  Chebotarev realizes, up to
conjugacy, every element of \(\Gamma\) as a Frobenius element.
Thus local vanishing almost everywhere is equivalent to vanishing
on every cyclic subgroup.  Conversely, a class with those cyclic
restrictions zero is locally zero at every unramified place, which
proves the identification.

Torus duality identifies the Hasse-principle obstruction for torsors
under \(T\) with the dual of \(\Sha^2(k,M)\), and controls the weak
approximation defect of \(T\) by \(\Sha_\omega^2(k,M)\)
\cite[Proposition 9.8 and Corollary 8.14]{Sansuc}.
The stated vanishing therefore gives
a rational point on every locally soluble torsor and weak
approximation on \(T\).  Translation by that rational point gives
weak approximation on the torsor.
\end{proof}

To calculate this kernel, it is convenient to use extensions by
lattices.  The following lemma will also apply when the finite
group is not a direct product.

\begin{lemma}[A central sign element]\label{lem:foundations-sign-extension}
Let \(\Gamma\) be a finite group and \(M\) a torsion-free
\(\Gamma\)-lattice.  Suppose that \(c\in Z(\Gamma)\) has order two
and acts on \(M\) as multiplication by \(-1\).
An extension
\[
       1\longrightarrow M\longrightarrow E
        \longrightarrow\Gamma\longrightarrow1
\]
determines a class
\(\beta(E)\in H^1(\Gamma/\langle c\rangle,M/2M)\), and this class
is zero if and only if the extension splits.
If the restriction to the subgroup generated by an involution
\(r\in\Gamma\) splits, the value at the image of \(r\) of every
cocycle representing \(\beta(E)\) lifts to an element
\(d_r\in M\) with \((1+r)d_r=0\).
\end{lemma}

\begin{proof}
Write the subgroup \(M\) additively, with translations on the left.
Let \(z\) lift \(c\).  Since \(z^2\in M\) commutes with \(z\),
it lies in \(M^c=0\); hence \(z\) is an involution.
Choose lifts \(s_g\) of \(g\in\Gamma\), with \(s_1=1\), and write
\[
 s_gs_h=f(g,h)s_{gh},\qquad
                 z s_gz^{-1}=d_gs_g.
\]
Comparing the two expressions for the conjugate of \(s_gs_h\) gives
\[
                    d_{gh}=d_g+g d_h+2f(g,h).
\]
Thus \(g\mapsto d_g\bmod 2M\) is a cocycle.
Replacing \(s_g\) by \(a_gs_g\) replaces \(d_g\) by
\(d_g-2a_g\).  Replacing \(z\) by \(az\) replaces \(d_g\) by
\(d_g+(1-g)a\), a coboundary change modulo two.
Choose \(s_c=z\); then \(d_c=0\), and the cocycle descends to
\(\Gamma/\langle c\rangle\), whose action on \(M/2M\) is well
defined.

If its class is zero, change \(z\) so that all \(d_g\) are even,
and then replace \(s_g\) by \((d_g/2)s_g\).  All these lifts commute
with \(z\).  The centralizer \(C_E(z)\) consequently maps onto
\(\Gamma\), and its kernel is \(M^c=0\).  The inverse of this
isomorphism is a splitting.  A splitting plainly gives the zero
class.  Finally, if \(s_r\) is an involutive lift of \(r\), conjugating
its square by \(z\) gives \((1+r)d_r=0\).
Changing the cocycle by a coboundary replaces this lift of its value
by \(d_r+(1-r)a\) for some \(a\in M\), and preserves the identity
because \((1+r)(1-r)a=0\).
\end{proof}

\subsection{The simultaneous norm torus}

\begin{lemma}[Hasse principle and approximation for a norm torus]
\label{lem:norm-torus}
Let \(L/k\) be quadratic, let \(F/k\) be separable of degree
\(m\geq3\), and suppose that the normal closure of \(F\), jointly
with \(L\), has Galois group \(S_m\times C_2\), with its natural
action on the embeddings of \(F\).  Set
\[
 T=\ker\left(
 \Res_{F/k}\bigl(\Res^1_{LF/F}\bbG_m\bigr)
       \xrightarrow{\,N_{LF/L}\,}\Res^1_{L/k}\bbG_m
 \right).
\]
This kernel is a torus.  Every torsor under \(T\) with points
over all completions has a rational point, and its rational points
are dense at any prescribed finite set of completions.
\end{lemma}

\begin{proof}
Let \(c\) denote the generator of the second factor \(C_2\).
On character lattices the norm map is the diagonal inclusion
\[
 \bbZ_{\mathrm{sign}}\longrightarrow
               \bbZ^m\otimes\mathrm{sign}_{C_2}.
\]
It is primitive.  Thus its cokernel is torsion-free, the kernel
of the norm map is connected, and its character lattice is
\begin{equation}\label{eq:foundations-norm-lattice}
       M=(\bbZ^m/\bbZ\mathbf1)\otimes\mathrm{sign}_{C_2}.
\end{equation}
Here \(S_m\) permutes the coordinates and \(c\) acts by \(-1\).
By Lemma~\ref{lem:foundations-cyclic-detection}, it suffices to split
every extension of \(S_m\times C_2\) by \(M\) that splits over
all cyclic subgroups.

Fix such an extension.  Lemma~\ref{lem:foundations-sign-extension}
gives a class
\[
                 \beta\in H^1(S_m,M/2M).
\]
We show that \(\beta=0\).  Put
\[
       V_m=\bbF_2^m,\qquad Q_m=V_m/\bbF_2\mathbf1=M/2M .
\]
The permutation-module sequence
\[
 0\longrightarrow\bbF_2\mathbf1
   \longrightarrow V_m\longrightarrow Q_m\longrightarrow0
\]
has a connecting map
\[
 H^1(S_m,Q_m)\longrightarrow H^2(S_m,\bbF_2).
\]
By Shapiro's lemma, the preceding map on degree-one cohomology
is the restriction
\[
 H^1(S_m,\bbF_2)\longrightarrow H^1(S_{m-1},\bbF_2).
\]
It is an isomorphism for \(m\geq3\): both sides are generated by
the permutation sign character.  Therefore the connecting map
is injective.  Its image consists of classes whose restriction to
the point stabilizer \(S_{m-1}\) is zero.

Let \(\widehat S_m\) be the central extension by \(\bbF_2\)
corresponding to the image of \(\beta\), and denote its nonidentity
central kernel element by \(\zeta\).  It splits on each point
stabilizer, by conjugacy.  Consequently every transposition has
involutive lifts.  For \(m\geq5\), two disjoint transpositions fix
a common point, so their lifts commute.  For \(m=3\) there is no
disjoint pair of Coxeter generators to consider.

In either of these cases, choose involutive lifts \(t_i\) of
the adjacent transpositions \((i,i+1)\).
Write
\[
                 (t_it_{i+1})^3=\zeta^{\epsilon_i},
                       \qquad \epsilon_i\in\bbF_2 .
\]
Multiplying \(t_i\) by \(\zeta^{a_i}\), with
\(a_1=0\) and \(a_{i+1}=a_i+\epsilon_i\), makes all these cubes
equal to \(1\).  It does not change the squares or the commutation
of distant generators.  The Coxeter presentation of \(S_m\)
therefore gives a splitting of \(\widehat S_m\).
By injectivity of the connecting map, \(\beta=0\).

It remains to treat \(m=4\), where disjoint transpositions have
no common fixed point.  The only remaining Coxeter relation is
commutation of lifts of
\(\tau=(12)\) and \(\nu=(34)\).
Fix the sign lift \(z\) in the proof of
Lemma~\ref{lem:foundations-sign-extension}, and use the cocycle
\(g\mapsto d_g\bmod2M\) produced there.  Choose representatives
\(b(g)\in\bbF_2^4\) of its values, with \(b(1)=0\).
Its connecting factor set is given by
\[
 \epsilon(g,h)\mathbf1=b(g)+g b(h)-b(gh).
\]
Since \(\tau\nu=\nu\tau\), failure of their lifts to commute would
give
\begin{equation}\label{eq:foundations-four-obstruction}
       (1+\nu)b(\tau)+(1+\tau)b(\nu)=\mathbf1 .
\end{equation}

We now return to the original extension by the lattice \(M\),
not the central extension \(\widehat S_4\).
Its restriction to \(\langle\tau\rangle\) splits by hypothesis.
Replace the chosen lift of \(\tau\) by an involutive lift.
This changes its translation parameter by twice an element of \(M\),
so the fixed cocycle modulo two is unchanged.  By
Lemma~\ref{lem:foundations-sign-extension}, the associated
translation parameter \(d\in M\) satisfies
\((1+\tau)d=0\), and its reduction represents \(\beta(\tau)\).
Choose an integral representative
\(\widetilde d\in\bbZ^4\).  For some \(a\in\bbZ\) we have
\[
           \widetilde d+\tau\widetilde d=a\mathbf1,
                \qquad 2\widetilde d_3=a=2\widetilde d_4.
\]
Hence \(\widetilde d_3=\widetilde d_4\).
Every representative \(b(\tau)\) has equal third and fourth
coordinates: adding a diagonal vector modulo two preserves that
equality.  The third and fourth coordinates of
\((1+\nu)b(\tau)\) are therefore zero, while those of
\((1+\tau)b(\nu)\) are zero because \(\tau\) fixes them.
This contradicts Equation~\ref{eq:foundations-four-obstruction}.
The disjoint lifts commute after all, and the same Coxeter-generator
adjustment splits \(\widehat S_4\).  Thus \(\beta=0\) also for \(m=4\).

In every degree \(m\geq3\),
Lemma~\ref{lem:foundations-sign-extension} now splits the original
lattice extension.  Lemma~\ref{lem:foundations-cyclic-detection}
gives \(\Sha_\omega^2(k,M)=0\) and the asserted Hasse principle
and weak approximation.
\end{proof}

We have obtained two distinct kinds of control.  The approximation
procedure chooses a representative inside a fixed abstract coset
while controlling its local centralizers.  The norm-torus lemma
then solves the simultaneous norm equations in those centralizers.
The next section uses this combination to construct commuting
rational elements with specified determinants.
